\subsection{Contextualización y análisis del problema}

\subsubsection{Conceptualización}

Un sistema de gestión de turnos es una herramienta tecnológica diseñada para organizar y administrar la atención de usuarios en un servicio, asignando turnos de manera ordenada y, en algunos casos, priorizada. En el contexto de los centros de salud de San Marcos, este tipo de sistema permite mejorar la distribución de pacientes, optimizar el tiempo de atención y garantizar un servicio más eficiente y equitativo.

El prototipo desarrollado se basa en la automatización del proceso de registro de pacientes, asignación de turnos y clasificación según nivel de prioridad (emergencia, adulto mayor, embarazada o atención normal). Su funcionamiento integra módulos que facilitan tanto la interacción del paciente como la gestión por parte del personal médico.

De esta manera, el sistema no solo organiza la atención, sino que también contribuye a mejorar la experiencia del usuario y el control operativo dentro del centro de salud.

\subsubsection{Análisis del problema}

En los centros de salud de (San Marcos Carazo), la gestión de turnos se realiza de forma manual o con sistemas poco eficientes, lo que genera diversas dificultades:

Largas filas de espera

Desorden en la atención de pacientes

Falta de priorización adecuada de casos urgentes

Pérdida de tiempo tanto para pacientes como para el personal

Sobrecarga en la organización del flujo de atención

Esta situación afecta directamente la calidad del servicio, provocando insatisfacción en los pacientes y dificultando el trabajo del personal médico.

El problema principal radica en la ausencia de un sistema automatizado que permita organizar de manera eficiente la atención, considerando criterios de prioridad y optimizando los recursos disponibles.

Por ello, surge la necesidad de desarrollar un prototipo de sistema de gestión de turnos que permita mejorar la organización, reducir los tiempos de espera y garantizar una atención más ordenada, rápida y eficiente.
